The reconstructed generator is a four-level 3D U-Net with channels $(32,64,64,128)$, one residual block per level, cross-attention only at the deepest level, and conditioning dimension 512. It contains 9,925,921 trainable parameters, excluding the frozen BiomedCLIP encoders. Generator training used Adam (learning rate $2.5\times10^{-5}$), batch size 1 and AMP. The principal generation endpoint is epoch 165; documented epoch-200 outputs and the MRI-only Swin-UNETR endpoint are distinguished from that endpoint. Recorded training/split settings by experiment family are summarized in Table~\ref{tab:reproducibility}; hardware and software differed between historical, reconstructed and timing runs. The historical estimate of 138 GPU-hours was not independently verified and is not used to claim full convergence or a speedup over other methods.

Inference timing used the Full fold-1 epoch-165 generator on an RTX 4090 (24 GB), PyTorch 2.5.1+cu121, CUDA FP32 and GPU-resident $1\times1\times160\times192\times160$ MRI with cached text context. After ten warm-up iterations, 50 synchronized repetitions per tracer gave FDG $251.015\pm0.187$ ms and AV45 $251.411\pm0.256$ ms (mean $\pm$ sample s.d. across repeated measurements of one fixed input). Peak PyTorch-allocated memory was 8.29 GiB and reserved memory 13.27 GiB, excluding other processes. These measurements exclude text encoding, preprocessing, transfers and file I/O; they describe a single-tracer forward pass, not end-to-end clinical workflow latency. Validated full-network FLOP counts and identical-protocol baseline timing were unavailable, so no FLOP estimate or baseline speedup ratio is claimed.
